\documentclass[12pt,a4paper]{article}

\usepackage[utf8]{inputenc}
\usepackage[T1]{fontenc}
\usepackage[spanish]{babel}
\usepackage{amsmath, amssymb, amsthm}
\usepackage{geometry}
\usepackage{booktabs}
\usepackage{listings}
\usepackage{xcolor}
\usepackage{hyperref}
\usepackage{parskip}

\geometry{margin=2.5cm}

\hypersetup{
    colorlinks=true,
    linkcolor=blue!60!black,
    urlcolor=blue!60!black,
    citecolor=blue!60!black
}

\lstset{
    language=Python,
    basicstyle=\ttfamily\small,
    keywordstyle=\color{blue!70!black}\bfseries,
    commentstyle=\color{green!50!black}\itshape,
    stringstyle=\color{orange!80!black},
    numbers=left,
    numberstyle=\tiny\color{gray},
    numbersep=8pt,
    frame=single,
    rulecolor=\color{gray!40},
    backgroundcolor=\color{gray!5},
    breaklines=true,
    showstringspaces=false,
    tabsize=4
}

\title{\textbf{Evaluación de Integrales de Dominio mediante\\
Dual Reciprocity y el Método de Soluciones Particulares}\\[0.5em]
\large Notas para un Curso de Maestría en Boundary Element Method}
\author{W.F. Florez\\
\small Wessex Institute of Technology}
\date{}

\begin{document}

\maketitle
\tableofcontents
\newpage

% ─────────────────────────────────────────────────────────────
\section{Planteamiento del Problema}
% ─────────────────────────────────────────────────────────────

El objetivo es puramente de cuadratura: convertir una integral de dominio

\begin{equation}
    I = \int_\Omega b(x)\, d\Omega_x
\end{equation}

en una suma de integrales de frontera, sin necesidad de discretizar el interior
del dominio $\Omega$. Esto se logra mediante el \textbf{método de soluciones
particulares} combinado con funciones de base radial (RBF), en el espíritu del
método de Dual Reciprocity (DR-BEM).

% ─────────────────────────────────────────────────────────────
\section{Conversión de Integral de Dominio a Integrales de Frontera}
% ─────────────────────────────────────────────────────────────

\subsection{Paso 1 — Aproximación RBF}

Se aproxima $b(x)$ en términos de funciones radiales evaluadas en $N$ nodos de
colocación $\{x_j\}_{j=1}^N \subset \Omega$:

\begin{equation}
    b(x) \approx \sum_{j=1}^{N} \alpha_j f_j(r_j), \qquad r_j = \|x - x_j\|
\end{equation}

\subsection{Paso 2 — Condición del Método de Soluciones Particulares}

Se exige que cada función base sea el Laplaciano de alguna función
$\hat{\phi}_j$:

\begin{equation}
    f_j(r_j) = \nabla^2 \hat{\phi}_j(r_j)
\end{equation}

de modo que la aproximación queda:

\begin{equation}
    b(x) \approx \sum_{j=1}^{N} \alpha_j \nabla^2 \hat{\phi}_j(r_j)
\end{equation}

\subsection{Paso 3 — Sustitución en la Integral}

\begin{equation}
    I \approx \sum_{j=1}^{N} \alpha_j \int_\Omega \nabla^2 \hat{\phi}_j\, d\Omega_x
\end{equation}

\subsection{Paso 4 — Teorema de la Divergencia}

Aplicando $\displaystyle\int_\Omega \nabla^2 \hat{\phi}\, d\Omega =
\int_\Gamma \nabla\hat{\phi} \cdot \mathbf{n}\, d\Gamma =
\int_\Gamma \frac{\partial \hat{\phi}}{\partial n}\, d\Gamma$:

\begin{equation}
\boxed{
    I \approx \sum_{j=1}^{N} \alpha_j \int_\Gamma
    \frac{\partial \hat{\phi}_j}{\partial n}\, d\Gamma_x
}
\end{equation}

La integral de dominio queda expresada como una suma de integrales de frontera,
cada una evaluable numéricamente con cuadratura estándar sobre $\Gamma$.

\subsection{Determinación de los Coeficientes $\alpha_j$ por Colocación}

Se conoce $b(x)$ en los $N$ nodos de colocación $\{x_i\}_{i=1}^N \subset
\Omega$. Evaluando la aproximación en cada punto:

\begin{equation}
    b(x_i) = \sum_{j=1}^{N} \alpha_j f_j(r_{ij}), \qquad
    r_{ij} = \|x_i - x_j\|
\end{equation}

lo que genera el sistema lineal:

\begin{equation}
    \mathbf{F}\boldsymbol{\alpha} = \mathbf{b}, \qquad F_{ij} = f_j(r_{ij})
\end{equation}

con solución:

\begin{equation}
    \boldsymbol{\alpha} = \mathbf{F}^{-1}\mathbf{b}
\end{equation}

\subsection{Generalización del Operador}

El operador no necesita ser $\nabla^2$. Sea $\mathcal{L}$ un operador
diferencial lineal para el cual existe una identidad de Green que permita
convertir integrales de dominio en integrales de frontera. Se exige:

\begin{equation}
    f_j = \mathcal{L}\hat{\phi}_j
\end{equation}

y la conversión queda:

\begin{equation}
    I \approx \sum_{j=1}^{N} \alpha_j \int_\Omega \mathcal{L}\hat{\phi}_j\,
    d\Omega_x = \sum_{j=1}^{N} \alpha_j \int_\Gamma
    \mathcal{B}[\hat{\phi}_j]\, d\Gamma_x
\end{equation}

donde $\mathcal{B}$ es el operador de frontera que surge de la identidad de
Green asociada a $\mathcal{L}$.

\textbf{Sobre la elección del operador:}
\begin{itemize}
    \item \textbf{Caso ideal:} $\mathcal{L}$ está relacionado con la naturaleza
    de $b(x)$, lo que garantiza que la aproximación sea eficiente y bien
    condicionada.
    \item \textbf{Caso general:} sin una PDE de fondo que sugiera $\mathcal{L}$,
    la elección es heurística. $\nabla^2$ es la opción natural por su
    simplicidad y porque las RBFs clásicas tienen $\hat{\phi}$ conocidas en
    forma cerrada.
    \item \textbf{Pregunta abierta pedagógicamente poderosa:} ¿cambia la
    calidad de la cuadratura resultante según el operador elegido,
    independientemente de que la integral de frontera sea exactamente evaluable?
\end{itemize}

% ─────────────────────────────────────────────────────────────
\section{Obtención de $\hat{\phi}(r)$ en 2D}
% ─────────────────────────────────────────────────────────────

En 2D se resuelve $\nabla^2\hat{\phi} = f(r)$ en coordenadas polares:

\begin{equation}
    \frac{1}{r}\frac{d}{dr}\left(r\frac{d\hat{\phi}}{dr}\right) = f(r)
\end{equation}

integrando dos veces directamente. La primera integración multiplica por $r$ e
integra; la constante de integración se toma cero por regularidad en $r=0$.

\subsection{Caso 1: $f(r) = 1 + r^3$}

\textbf{Primera integración:}

\begin{equation}
    \frac{d}{dr}\left(r\frac{d\hat{\phi}}{dr}\right) = r + r^4
    \implies r\frac{d\hat{\phi}}{dr} = \frac{r^2}{2} + \frac{r^5}{5}
    \implies \frac{d\hat{\phi}}{dr} = \frac{r}{2} + \frac{r^4}{5}
\end{equation}

\textbf{Segunda integración:}

\begin{equation}
\boxed{
    \hat{\phi}(r) = \frac{r^2}{4} + \frac{r^5}{25}
}
\end{equation}

Completamente algebraica, sin logaritmos ni funciones especiales.

\subsection{Caso 2: Multiquadrics $f(r) = (r^2 + c^2)^{1/2}$}

\textbf{Primera integración:}

\begin{equation}
    r\frac{d\hat{\phi}}{dr} = \int r(r^2+c^2)^{1/2}\,dr
    = \frac{1}{3}(r^2+c^2)^{3/2}
    \implies \frac{d\hat{\phi}}{dr} = \frac{(r^2+c^2)^{3/2}}{3r}
\end{equation}

\textbf{Segunda integración} (usando la integral estándar
$\int \frac{(r^2+c^2)^{3/2}}{r}\,dr$):

\begin{equation}
\boxed{
    \hat{\phi}(r) = \frac{(r^2+c^2)^{3/2}}{9}
    + \frac{c^2\sqrt{r^2+c^2}}{3}
    - \frac{c^3}{3}\ln\!\left(\frac{c+\sqrt{r^2+c^2}}{r}\right)
}
\end{equation}

El término logarítmico no diverge en $r=0$ cuando $c>0$, ya que el parámetro
de forma $c$ regulariza la singularidad. Esta es una ventaja importante de la
multiquadric sobre otras RBFs.

\subsection{Tabla Comparativa}

\begin{center}
\begin{tabular}{lll}
\toprule
$f(r)$ & $\hat{\phi}(r)$ & Observación \\
\midrule
$1 + r^3$ & $\dfrac{r^2}{4} + \dfrac{r^5}{25}$ & Completamente algebraica \\[1.2em]
$(r^2+c^2)^{1/2}$ &
$\dfrac{(r^2+c^2)^{3/2}}{9} + \dfrac{c^2\sqrt{r^2+c^2}}{3}
- \dfrac{c^3}{3}\ln\!\left(\dfrac{c+\sqrt{r^2+c^2}}{r}\right)$ &
Logaritmo regularizado por $c>0$ \\
\bottomrule
\end{tabular}
\end{center}

\textbf{Observación:} para $f(r) = \sum a_k r^k$ (polinomio en $r$), la
integración produce siempre potencias de $r$ — forma cerrada algebraica
garantizada. La multiquadric introduce el término logarítmico inevitablemente
por la raíz cuadrada.

% ─────────────────────────────────────────────────────────────
\section{Derivada Normal $\partial\hat{\phi}/\partial n$}
% ─────────────────────────────────────────────────────────────

Usando la regla de la cadena, dado que $\hat{\phi} = \hat{\phi}(r)$ con
$r = \|x - x_j\|$:

\begin{equation}
    \frac{\partial\hat{\phi}}{\partial n}
    = \frac{d\hat{\phi}}{dr}\frac{\partial r}{\partial n}
    = \frac{d\hat{\phi}}{dr}\,\frac{(x - x_j)\cdot\mathbf{n}}{r}
\end{equation}

El factor $\frac{(x-x_j)\cdot\mathbf{n}}{r} = \cos\theta$ es puramente
geométrico y común a ambos casos.

\subsection{Caso 1: $f(r) = 1 + r^3$}

\begin{equation}
    \frac{d\hat{\phi}}{dr} = \frac{r}{2} + \frac{r^4}{5}
\end{equation}

\begin{equation}
\boxed{
    \frac{\partial\hat{\phi}}{\partial n}
    = \left(\frac{1}{2} + \frac{r^3}{5}\right)(x - x_j)\cdot\mathbf{n}
}
\end{equation}

Completamente regular para cualquier configuración de nodos.

\subsection{Caso 2: Multiquadrics $f(r) = (r^2+c^2)^{1/2}$}

Derivando $\hat{\phi}$ término a término:

\begin{align}
    \frac{d}{dr}\left[\frac{(r^2+c^2)^{3/2}}{9}\right]
    &= \frac{r(r^2+c^2)^{1/2}}{3} \\[0.5em]
    \frac{d}{dr}\left[\frac{c^2(r^2+c^2)^{1/2}}{3}\right]
    &= \frac{c^2 r}{3(r^2+c^2)^{1/2}} \\[0.5em]
    \frac{d}{dr}\left[-\frac{c^3}{3}
    \ln\!\left(\frac{c+\sqrt{r^2+c^2}}{r}\right)\right]
    &= \frac{c^3}{3r\sqrt{r^2+c^2}}
\end{align}

Sumando y simplificando:

\begin{equation}
    \frac{d\hat{\phi}}{dr} = \frac{r^4 + c^2 r^2 + c^3}{3r\sqrt{r^2+c^2}}
\end{equation}

\begin{equation}
\boxed{
    \frac{\partial\hat{\phi}}{\partial n}
    = \frac{r^4 + c^2 r^2 + c^3}{3r^2\sqrt{r^2+c^2}}
    \,(x - x_j)\cdot\mathbf{n}
}
\end{equation}

\textbf{Regularidad:} cuando $c>0$, al tomar $r\to 0$ se tiene
$\frac{\partial\hat{\phi}}{\partial n} \to \frac{c^2}{3}\,\frac{(x-x_j)\cdot\mathbf{n}}{r}$,
donde el factor $\cos\theta$ es el que controla el comportamiento near-singular —
exactamente como en el kernel BEM clásico $\partial u^*/\partial n$.

% ─────────────────────────────────────────────────────────────
\section{Análisis de Singularidades}
% ─────────────────────────────────────────────────────────────

\textbf{Configuración geométrica:} los nodos de colocación $\{x_j\}$ están en
$\Omega$ (interior), y la integración $\int_\Gamma \frac{\partial\hat{\phi}_j}
{\partial n}\,d\Gamma$ involucra el radio $r_j = \|x - x_j\|$ con $x\in\Gamma$.
Dado que $x_j\in\Omega$ y $x\in\Gamma$:

\begin{equation}
    r_{j,\min} = \mathrm{dist}(x_j, \Gamma) > 0
\end{equation}

por lo que \textbf{no hay singularidad verdadera} en la integración sobre
$\Gamma$ — se trata de una \textbf{near-singularity} cuya severidad depende de
qué tan cerca estén los nodos internos de la frontera.

\textbf{El caso problemático} surge si se permite $x_j\in\Gamma$ (colocación
sobre la frontera). En ese caso $r_j\to 0$ cuando el punto de integración
$x\in\Gamma$ se acerca a $x_j\in\Gamma$, produciendo:

\begin{itemize}
    \item Para $1+r^3$: $\hat{\phi}$ regular, pero $d\hat{\phi}/dr\to 0$
    conforme $r\to 0$, sin singularidad.
    \item Para MQ: singularidad logarítmica en $\frac{\partial\hat{\phi}}
    {\partial n}$ cuando $r\to 0$.
\end{itemize}

Las estrategias disponibles para ese caso son las del BEM clásico:
integración analítica local, sustracción de singularidad, o transformaciones
de coordenadas (transformación de Telles).

% ─────────────────────────────────────────────────────────────
\section{Vector Normal sobre Curvas y Superficies Paramétricas}
% ─────────────────────────────────────────────────────────────

\subsection{Caso 2D: Curva Plana Parametrizada}

Sea la curva $\Gamma$ parametrizada por $t\in[a,b]$:
$\mathbf{x}(t) = (x(t),\, y(t))^T$.

El \textbf{vector tangente} es $\mathbf{T} = \mathbf{x}'(t) = (x'(t),\,y'(t))^T$.

El \textbf{vector normal exterior unitario} se obtiene rotando $\mathbf{T}$
en $90°$ hacia afuera:

\begin{equation}
    \mathbf{n} = \frac{1}{\|\mathbf{T}\|}
    \begin{pmatrix} y'(t) \\ -x'(t) \end{pmatrix},
    \qquad \|\mathbf{T}\| = \sqrt{x'^2 + y'^2} = \frac{d\Gamma}{dt}
\end{equation}

La derivada normal integrada sobre $\Gamma$ queda:

\begin{equation}
    \frac{\partial\hat{\phi}}{\partial n}\,d\Gamma
    = \nabla\hat{\phi}\cdot\mathbf{n}\,d\Gamma
    = \left(\frac{\partial\hat{\phi}}{\partial x}y'
    - \frac{\partial\hat{\phi}}{\partial y}x'\right)dt
\end{equation}

El jacobiano \textbf{cancela} al normalizador — no hace falta calcular
$\|\mathbf{T}\|$ explícitamente.

\subsection{Caso 3D: Superficie Parametrizada}

Sea la superficie $\Gamma$ parametrizada por $(s,t)\in\mathcal{D}$:
$\mathbf{x}(s,t) = (x,y,z)^T$.

Los \textbf{vectores tangentes} son $\mathbf{T}_s = \partial\mathbf{x}/\partial s$
y $\mathbf{T}_t = \partial\mathbf{x}/\partial t$.

El \textbf{vector normal} (no unitario):

\begin{equation}
    \mathbf{N} = \mathbf{T}_s\times\mathbf{T}_t
    = \begin{vmatrix}
        \mathbf{e}_1 & \mathbf{e}_2 & \mathbf{e}_3 \\
        \partial x/\partial s & \partial y/\partial s & \partial z/\partial s \\
        \partial x/\partial t & \partial y/\partial t & \partial z/\partial t
    \end{vmatrix}
\end{equation}

El \textbf{elemento de área}: $d\Gamma = \|\mathbf{N}\|\,ds\,dt$.

La derivada normal integrada:

\begin{equation}
    \frac{\partial\hat{\phi}}{\partial n}\,d\Gamma
    = \nabla\hat{\phi}\cdot\mathbf{N}\,ds\,dt
\end{equation}

El normalizador cancela de nuevo — solo se necesita $\mathbf{N}$ sin normalizar.

\subsection{Tabla Resumen}

\begin{center}
\begin{tabular}{lll}
\toprule
 & 2D & 3D \\
\midrule
Parámetros & $t$ & $s,\,t$ \\
Tangentes & $\mathbf{T} = \mathbf{x}'(t)$ & $\mathbf{T}_s,\,\mathbf{T}_t$ \\
Normal no unitaria & $(y',\,-x')$ & $\mathbf{T}_s\times\mathbf{T}_t$ \\
Jacobiano & $\sqrt{x'^2+y'^2}$ & $\|\mathbf{T}_s\times\mathbf{T}_t\|$ \\
Forma integrada & $\nabla\hat{\phi}\cdot(y',-x')\,dt$
    & $\nabla\hat{\phi}\cdot\mathbf{N}\,ds\,dt$ \\
\bottomrule
\end{tabular}
\end{center}

% ─────────────────────────────────────────────────────────────
\section{Discretización con Elementos Constantes en 2D}
% ─────────────────────────────────────────────────────────────

\subsection{Geometría del Elemento}

Se divide $\Gamma$ en $N_e$ elementos rectos $\Gamma_k$, cada uno definido por
dos nodos extremos $\mathbf{x}_k^1$ y $\mathbf{x}_k^2$, con el
\textbf{nodo de colocación central}:

\begin{equation}
    \mathbf{x}_k^c = \frac{\mathbf{x}_k^1 + \mathbf{x}_k^2}{2}
\end{equation}

Parametrización local con $t\in[-1,1]$:

\begin{equation}
    \mathbf{x}(t) = \frac{1-t}{2}\mathbf{x}_k^1 + \frac{1+t}{2}\mathbf{x}_k^2
\end{equation}

Las derivadas son constantes sobre cada elemento:

\begin{equation}
    x'(t) = \frac{x_k^2 - x_k^1}{2}, \qquad y'(t) = \frac{y_k^2 - y_k^1}{2}
\end{equation}

\subsection{Jacobiano y Normal}

El jacobiano es constante en cada elemento:

\begin{equation}
    J_k = \frac{L_k}{2},
    \qquad L_k = \sqrt{(x_k^2-x_k^1)^2 + (y_k^2-y_k^1)^2}
\end{equation}

El vector normal exterior unitario, también constante:

\begin{equation}
    \mathbf{n}_k = \frac{1}{L_k}
    \begin{pmatrix} y_k^2 - y_k^1 \\ -(x_k^2 - x_k^1) \end{pmatrix}
\end{equation}

\subsection{La Integral sobre $\Gamma$}

\begin{equation}
    \int_\Gamma \frac{\partial\hat{\phi}_j}{\partial n}\,d\Gamma
    = \sum_{k=1}^{N_e} \int_{\Gamma_k} \frac{\partial\hat{\phi}_j}{\partial n}
    \,d\Gamma
\end{equation}

Con cuadratura de Gauss de $n_g$ puntos $\{t_p, w_p\}$:

\begin{equation}
    \int_{\Gamma_k}\frac{\partial\hat{\phi}_j}{\partial n}\,d\Gamma
    \approx J_k \sum_{p=1}^{n_g} w_p\,\frac{d\hat{\phi}}{dr}\bigg|_{r_{jp}}
    \frac{(\mathbf{x}(t_p) - \mathbf{x}_j)\cdot\mathbf{n}_k}{r_{jp}}
\end{equation}

\subsection{Resultado Final Ensamblado}

\begin{equation}
    \int_\Omega b(x)\,d\Omega \approx
    \sum_{j=1}^{N}\alpha_j \sum_{k=1}^{N_e} J_k
    \sum_{p=1}^{n_g} w_p
    \frac{d\hat{\phi}}{dr}\bigg|_{r_{jp}}
    \frac{(\mathbf{x}(t_p)-\mathbf{x}_j)\cdot\mathbf{n}_k}{r_{jp}}
\end{equation}

Matricialmente, definiendo el vector $\mathbf{q}$ con componentes:

\begin{equation}
    q_j = \sum_{k=1}^{N_e} J_k \sum_{p=1}^{n_g} w_p
    \frac{d\hat{\phi}}{dr}\bigg|_{r_{jp}}
    \frac{(\mathbf{x}(t_p)-\mathbf{x}_j)\cdot\mathbf{n}_k}{r_{jp}}
\end{equation}

el resultado final es simplemente:

\begin{equation}
    \boxed{I \approx \mathbf{q}^T\boldsymbol{\alpha}
    = \mathbf{q}^T \mathbf{F}^{-1}\mathbf{b}}
\end{equation}

% ─────────────────────────────────────────────────────────────
\section{Pseudocódigo}
% ─────────────────────────────────────────────────────────────

\begin{lstlisting}[language={},basicstyle=\ttfamily\small,
                   frame=single,backgroundcolor=\color{gray!5},
                   rulecolor=\color{gray!40}]
DATOS GEOMETRICOS
  Definir 4 elementos sobre los 4 lados del cuadrado [0,1]^2
  Cada elemento: nodo extremo 1, nodo extremo 2, nodo central, normal exterior

NODOS DE COLOCACION INTERNOS
  Colocar N puntos x_j dentro de Omega (no sobre Gamma)
  Evaluar b(x_j) en cada nodo -> vector b

SISTEMA DE COLOCACION RBF
  Para cada par (i,j) calcular r_ij = ||x_i - x_j||
  Armar matriz F con F_ij = f(r_ij) = 1 + r_ij^3
  Resolver F * alpha = b  ->  alpha = F^{-1} * b

VECTOR q  (integrales de frontera)
  Para cada nodo fuente j = 1,...,N
    q_j = 0
    Para cada elemento k = 1,...,4
      Para cada punto de Gauss p = 1,...,ng
        Calcular x(tp) sobre elemento k
        Calcular r_jp = ||x(tp) - x_j||
        Calcular dphi_dr = r_jp/2 + r_jp^4/5
        Calcular (x(tp) - x_j) . n_k
        q_j += w_p * J_k * (dphi_dr / r_jp) * dot

RESULTADO
  I_aprox = q^T * alpha
  Comparar con I_analitica = 2/3
\end{lstlisting}

% ─────────────────────────────────────────────────────────────
\section{Ejemplo Numérico: Dominio Cuadrado $[0,1]^2$}
% ─────────────────────────────────────────────────────────────

\subsection{Configuración}

\begin{itemize}
    \item Dominio: $\Omega = [0,1]^2$
    \item Función: $b(x,y) = x^2 + y^2$
    \item RBF: $f(r) = 1 + r^3 \implies \hat{\phi}(r) = r^2/4 + r^5/25$
    \item Elementos de contorno: 4 elementos constantes, uno por cada lado
    \item Nodos de colocación internos: grilla $4\times 4 = 16$ puntos en
    $[0.1, 0.9]^2$
    \item Puntos de Gauss por elemento: $n_g = 4$
\end{itemize}

\textbf{Solución analítica:}

\begin{equation}
    I = \int_0^1\int_0^1 (x^2+y^2)\,dx\,dy
    = \frac{1}{3} + \frac{1}{3} = \frac{2}{3}
\end{equation}

La geometría de los 4 elementos es:

\begin{center}
\begin{tabular}{llll}
\toprule
Lado & Extremo 1 & Extremo 2 & Normal $\mathbf{n}_k$ \\
\midrule
Bottom & $(0,0)$ & $(1,0)$ & $(0,-1)$ \\
Right  & $(1,0)$ & $(1,1)$ & $(1,0)$  \\
Top    & $(1,1)$ & $(0,1)$ & $(0,1)$  \\
Left   & $(0,1)$ & $(0,0)$ & $(-1,0)$ \\
\bottomrule
\end{tabular}
\end{center} 

\subsection{Código Python}

\begin{lstlisting}[language=Python]
"""
Evaluacion de integral de dominio mediante Dual Reciprocity (DR)
Dominio: cuadrado unitario [0,1]^2
Funcion: b(x,y) = x^2 + y^2
RBF: f(r) = 1 + r^3  =>  phi(r) = r^2/4 + r^5/25
Elementos constantes: 4 elementos, uno por cada lado
Solucion analitica: I = 2/3
"""

import numpy as np


# 1. FUNCIONES RBF Y SOLUCION PARTICULAR

def f_rbf(r):
    """RBF: f(r) = 1 + r^3"""
    return 1.0 + r**3

def phi(r):
    """Solucion particular: phi(r) = r^2/4 + r^5/25"""
    return r**2 / 4.0 + r**5 / 25.0

def dphi_dr(r):
    """Derivada: dphi/dr = r/2 + r^4/5"""
    return r / 2.0 + r**4 / 5.0

def b_func(x, y):
    """Funcion a integrar"""
    return x**2 + y**2

# 2. GEOMETRIA: 4 ELEMENTOS CONSTANTES

elements = np.array([
    [0.0, 0.0,  1.0, 0.0],   # bottom
    [1.0, 0.0,  1.0, 1.0],   # right
    [1.0, 1.0,  0.0, 1.0],   # top
    [0.0, 1.0,  0.0, 0.0],   # left
])

normals = np.array([
    [ 0.0, -1.0],
    [ 1.0,  0.0],
    [ 0.0,  1.0],
    [-1.0,  0.0],
])

def get_jacobian(elem):
    x1, y1, x2, y2 = elem
    L = np.sqrt((x2-x1)**2 + (y2-y1)**2)
    return L / 2.0

def get_point(elem, t):
    x1, y1, x2, y2 = elem
    x = 0.5*(1-t)*x1 + 0.5*(1+t)*x2
    y = 0.5*(1-t)*y1 + 0.5*(1+t)*y2
    return np.array([x, y])

# 3. PUNTOS DE GAUSS en [-1,1]

n_gauss = 4
gauss_pts, gauss_wts = np.polynomial.legendre.leggauss(n_gauss)

# 4. NODOS DE COLOCACION INTERNOS

n_side = 4
coords = np.linspace(0.1, 0.9, n_side)
collocation_nodes = []
for xi in coords:
    for yi in coords:
        collocation_nodes.append([xi, yi])
collocation_nodes = np.array(collocation_nodes)
N = len(collocation_nodes)

# 5. SISTEMA RBF: F * alpha = b

F_mat = np.zeros((N, N))
for i in range(N):
    for j in range(N):
        r = np.linalg.norm(collocation_nodes[i] - collocation_nodes[j])
        F_mat[i, j] = f_rbf(r)

b_vec = np.array([b_func(p[0], p[1]) for p in collocation_nodes])
alpha = np.linalg.solve(F_mat, b_vec)

 6. VECTOR q: integrales de frontera

q_vec = np.zeros(N)

for j in range(N):
    xj = collocation_nodes[j]
    q_j = 0.0
    for k in range(4):
        elem = elements[k]
        nk   = normals[k]
        Jk   = get_jacobian(elem)
        for p in range(n_gauss):
            tp   = gauss_pts[p]
            wp   = gauss_wts[p]
            xp   = get_point(elem, tp)
            rvec = xp - xj
            r    = np.linalg.norm(rvec)
            if r < 1e-14:
                continue
            dphi = dphi_dr(r)
            dot  = np.dot(rvec, nk)
            q_j += wp * Jk * (dphi / r) * dot
    q_vec[j] = q_j

# 7. RESULTADO

I_dr        = np.dot(q_vec, alpha)
I_analitica = 2.0 / 3.0
error       = abs(I_dr - I_analitica) / I_analitica * 100

print(f"Integral analitica     : {I_analitica:.10f}")
print(f"Integral DR (numerica) : {I_dr:.10f}")
print(f"Error relativo         : {error:.6f} %")
\end{lstlisting}

\subsection{Resultado}

\begin{center}
\begin{tabular}{ll}
\toprule
Cantidad & Valor \\
\midrule
Integral analítica      & $0.6666666667$ \\
Integral DR (numérica)  & $0.6663737442$ \\
Error relativo          & $0.043938\%$   \\
\bottomrule
\end{tabular}
\end{center}

Con solo 16 nodos internos, 4 elementos constantes y 4 puntos de Gauss,
el error es menor al $0.05\%$.

% ─────────────────────────────────────────────────────────────
\section{Observaciones Pedagógicas y Extensiones}
% ─────────────────────────────────────────────────────────────

\subsection{Refinamiento y Convergencia}

Hay tres experimentos naturales para los estudiantes:

\begin{enumerate}
    \item \textbf{Aumentar $n\_side$} de 4 a 6 u 8 y observar la caída del
    error. Esto ilustra la convergencia del sistema de colocación RBF, no de
    la cuadratura de Gauss, que ya es exacta con 4 puntos para funciones
    suaves.
    \item \textbf{Más elementos de contorno}: pasar de 4 a 8 o 16 elementos
    por lado. El efecto es menor porque con elementos constantes la geometría
    del cuadrado ya es exacta.
    \item \textbf{Condicionamiento de $\mathbf{F}$}: imprimir
    \texttt{np.linalg.cond(F\_mat)} para distintas distribuciones de nodos,
    introduciendo la discusión de estabilidad de la interpolación RBF.
\end{enumerate}

\subsection{Pregunta de Examen Sugerida}

\textit{¿Para qué clase de $f(r)$ está garantizada una $\hat{\phi}$ algebraica
en 2D?} La respuesta es inmediata desde la estructura de la integración en
coordenadas polares: para todo $f$ que sea un polinomio en $r$.

\subsection{Limitaciones Conocidas}

\begin{itemize}
    \item No se pueden usar nodos de colocación sobre la frontera sin
    técnicas de evaluación de integrales singulares o remoción de
    singularidades.
    \item La matriz $\mathbf{F}$ puede ser mal condicionada si los nodos
    están mal distribuidos — punto de entrada natural a la discusión de
    condicionamiento numérico en interpolación RBF.
    \item La calidad de la cuadratura depende enteramente de qué tan bien
    $\{f_j\}$ aproxima $b(x)$ — conexión directa con teoría de aproximación.
\end{itemize}

\end{document}
